\subsection{乘法群 U}

定理 1. 绝对值为 1 的复数的（乘法）拓扑群 U 同构于模 1 实数的（加法）拓扑群 T。

\(\mathbf{U} = \mathbf{S}_{1}\) 是紧且连通的，并且它的单位元 \(+1\) 有一个同胚于 \(\mathbf{R}\) 的开区间的邻域（第六章，§ 2，第 4 号，命题 5）；因此该定理是第五章 § 3（定理 2）所给的 \(\mathbf{T}\) 的拓扑刻画的一个推论。

推论. 非零复数的乘法群 \(\mathbf{C}^*\) 同构于群 \(\mathbf{R} \times \mathbf{T}\) （参看 § 1. 第 3 号，命题 1）。

注. 群 \(C^{*}\) 与 \(R \times T\) 的同构蕴含着在域 C 中每个“二项方程” \(z^{n} = a\) 都有根。利用这一事实以及 C 的局部紧性，我们可以得到达朗贝尔–高斯定理的另一个证明（习题 2）。

从群 \(\mathbf{T}\) 到群 \(\mathbf{U}\) 上的不同同构只有两个；因为若 \(g\) 、\(g'\) 是 \(\mathbf{T}\) 到 \(\mathbf{U}\) 上的两个同构，而 \(h'\) 是 \(g'\) 的逆，则 \(h' \circ g\) 是 \(\mathbf{T}\) 的一个自同构，因而（第七章，§ 2，第 4 号，命题 6）我们恒有 \(g'(x) = g(x)\) 或 \(g'(x) = g(-x)\) 。我们总可以假设 \(g\) 使得 \(i\) 是 \(g\) 对模 \(1\) 类（点 \(\frac{1}{4}\) 的）所取的像；于是，若 \(\varphi\) 表示 \(\mathbf{R}\) 到 \(\mathbf{T}\) 上的典范同态，则加法群 \(\mathbf{R}\) 到乘法群 \(\mathbf{U}\) 上的每个严格态射都形如 \(x \to g(\varphi(x/a))\) ，其中 \(a\) 是一个 \(\neq o\) 的实数（第七章，§ 2，第 3 号，命题 4）；注意，区间 \([- \frac{1}{2} |a|, \frac{1}{2} |a|]\) 是 \(\mathbf{R}\) 中被这个严格态射一一地映到其像上的最大对称开区间，并且我们有 \(g(\varphi(\frac{1}{4})) = i\) 。我们将把同态 \(x \to g(\varphi(x))\) 记作 \(x \to e(x)\) ；因而 \(\mathbf{R}\) 到 \(\mathbf{U}\) 上的每个严格态射都形如 \(x \to e(x/a)\) ，其中 \(a \neq 0\) 。函数 \(e(x)\) 在 \(\mathbf{R}\) 上连续、取复值，并且满足恒等式

\begin{enumerate}
\def\labelenumi{(\arabic{enumi})}
\tightlist
\item
\end{enumerate}

\[
\left| e (x) \right| = 1,\tag{2}
\]

\[
\boldsymbol {e} (x + y) = \boldsymbol {e} (x) \boldsymbol {e} (y),
\]

以及关系式

\[
e (0) = \mathrm{I}, \quad e (\frac {1}{4}) = i, \quad e (\frac {1}{2}) = - \mathrm{I}, \quad e (\frac {3}{4}) = - i, \quad e (\mathrm{I}) = \mathrm{I}. \tag {3}
\]

由 (1) 与 (2) 得出

\[
e (- x) = \frac {1}{e (x)} = \overline {{{{e (x)}}}},\tag{4}
\]

而由 (2) 与 (3) 得出

\[
\begin{array}{l} e (x + \frac {1}{4}) = i e (x), \quad e (x + \frac {1}{2}) = - e (x), \\ e (x + \frac {3}{4}) = - i e (x), \quad e (x + 1) = e (x). \end{array}
\]

于是函数 \(e(x)\) 是周期的，并且以 1 为基本周期。

注. 映射 \(x + iy \to e^x e(y)\) 是加法群 \(\mathbf{C}\) 到乘法群 \(\mathbf{C}^*\) 上的一个严格态射，并且它在 \(\mathbf{o}\) 的一个适当邻域上的限制是 \(\mathbf{C}\) 与 \(\mathbf{C}^*\) 的一个局部同构。因而（第七章，§ 2，第 3 号）\(\mathbf{C}\) 到 \(\mathbf{C}^*\) 上的每个严格态射都形如 \(x + iy \to e^{\alpha x + \beta y} e(\gamma x + \delta y)\) ，其中 \(\alpha, \beta, \gamma, \delta\) 是使 \(\alpha \delta - \beta \gamma \neq \mathbf{o}\) 的任意实数。我们后面将看到，这些同态中恰有一个，记作 \(z \rightarrow e^{z}\) ，使得

\[
\lim _ {z \rightarrow 0} \frac {e ^ {z} - 1}{z} = 1;
\]

并且这个同态在实轴上的限制与 \(e^{x}\) 相同（由此得到该记号）。

